\subsection{Hilbert-Samuel series of a well-filtered module}

In the remainder of this paragraph, we will use the following notation: if \(\mathbf{G} \in \mathbf{Z}((\mathbf{T}))\) and if \(r \in \mathbf{N}\) , we set \(\mathbf{G}^{(r)} = (1 - \mathbf{T})^{-r}\mathbf{G}\) ; in particular, if \(\mathbf{G} = \sum_{n \in \mathbf{Z}} a_n \mathbf{T}^n\) , then

\[
\mathrm{G} ^ {(1)} = \sum_ {n \in \mathbf {Z}} \left(\sum_ {i \leqslant n} a _ {i}\right) \mathrm{T} ^ {n}.
\]

If \(\mathbf{G} \geqslant 0\) , one has \(\mathbf{G}^{(r)} \geqslant 0\) for every \(r \in \mathbf{N}\) .

Let A be a Noetherian ring, q an ideal of A, and M a finitely generated A-module. Recall (III, § 3, no. 1, def. 1) that a q-good filtration on M is a map F : n ↦ Fn from Z to the set of submodules of M satisfying the following three conditions:

\begin{enumerate}
\def\labelenumi{\alph{enumi})}
\item
  we have \(\mathfrak{q}F_n\subset \mathrm{F}_{n + 1}\subset \mathrm{F}_n\) for every \(n\in \mathbf{Z},\)
\item
  there exists \(n_0 \in \mathbf{Z}\) such that \(\mathfrak{q}F_n = \mathrm{F}_{n+1}\) for \(n \geqslant n_0\) ,
\item
  there exists \(n_{1} \in \mathbf{Z}\) such that \(\mathrm{F}_{n_1} = \mathbf{M}\) .
\end{enumerate}

If \(n_{0}\) and \(n_{1}\) satisfy the preceding conditions, we have, for all \(n\in\mathbf{Z}\) ,

\[
\mathfrak {q} ^ {n - n _ {1}} \mathrm{M} \subset \mathrm{F} _ {n} \subset \mathfrak {q} ^ {n - n _ {0}} \mathrm{M}\tag{7}
\]

(recall that we have set \(\mathfrak{q}^r = \mathbf{A}\) for \(r\leqslant 0\) , by convention).

Lemma 3.--- If F and F' are two q-good filtrations on M, there exists an integer m such that \(F_{n}^{\prime} \subset F_{n-m}\) for every \(n \in Z\) .

Indeed, there exists \(n_{2}\) such that \(\mathbf{F}_{n}^{\prime}\subset\mathfrak{q}^{n-n_{2}}\mathbf{M}\) for every \(n\) , hence \(\mathbf{F}_{n}^{\prime}\subset\mathbf{F}_{n-(n_{2}-n_{1})}\) for every \(n\) .

Lemma 4. --- Let F be a q-good filtration on M. If M/qM has finite length, then M/F \(_{n+1}\) and F \(_{n}\) /F \(_{n+1}\) have finite length for all n ∈ Z.

With the notation from (7), we have long(M/F \(_{n+1}\) ) ≤ long(M/q \(^{n-n_1+1}\) M), and it suffices to prove that q \(^{n}\) M/q \(^{n+1}\) M has finite length for every n. We are therefore reduced to the case of the q-adic filtration. Let (x \(_{1}\) , \ldots, x \(_{r}\) ) be a finite generating system of the A-module q, and let I be the finite set of monomials of total degree n in r variables

\(X_{1},...,X_{r}\) . The homomorphism from \((\mathbf{M}/\mathfrak{q}\mathbf{M})^{\mathbf{I}}\) into \(q^{n}M/q^{n+1}M\) which maps a family \((u_{m})_{m\in\mathbf{I}}\) to the element \(\sum m(x_{1},...,x_{r})u_{m}\) is surjective. Since M/qM has finite length, \(q^{n}M/q^{n+1}M\) also has finite length.

Assume now that M/qM has finite length. Let F be a q-good filtration on M. There exists \(n_{1} \in Z\) such that \(F_{n_{1}} = M\) , hence \(F_{n} = M\) for \(n \leqslant n_{1}\) ; we therefore define an element \(H_{M,F}\) of \(\mathbf{Z}((T))\) by setting

\[
\mathrm{H} _ {\mathrm{M}, \mathrm{F}} = \sum_ {n \in \mathbf {Z}} \operatorname{long} _ {\mathrm{A} / \mathfrak {q}} (\mathrm{F} _ {n} / \mathrm{F} _ {n + 1}). \mathrm{T} ^ {n} \in \mathbf {Z} ((\mathrm{T})) .\tag{8}
\]

DEFINITION 1. --- We call H \(_{M,F}\) the Hilbert-Samuel series of the A-module M (relative to the q-good filtration F).

The map \(n\mapsto\text{long}_{\mathbf{A}}(\mathbf{F}_{n}/\mathbf{F}_{n+1})\) is often called the Hilbert-Samuel function of M (relative to F).

This applies in particular to the case of the q-adic filtration (\(F_{n} = q^{n}M\)); we then set \(H_{M,F} = H_{M,q}\) . One therefore has

\[
\mathrm{H} _ {\mathrm{M}, \mathfrak {q}} = \sum_ {n \in \mathbf {Z}} \operatorname{long} _ {\mathrm{A} / \mathfrak {q}} (\mathfrak {q} ^ {n} \mathrm{M} / \mathfrak {q} ^ {n + 1} \mathrm{M}). \mathrm{T} ^ {n}.\tag{9}
\]

PROPOSITION 3.--- a) If F is a q-good filtration on M, then we have

\[
\mathrm{H} _ {\mathrm{M}, \mathrm{F}} ^ {(1)} = \sum_ {n \in \mathbf {Z}} \operatorname{long} _ {\mathrm{A}} (\mathrm{M} / \mathrm{F} _ {n + 1}). \mathrm{T} ^ {n}.\tag{10}
\]

\begin{enumerate}
\def\labelenumi{\alph{enumi})}
\setcounter{enumi}{1}
\tightlist
\item
  If F and F' are two q-good filtrations on M, there exists an integer m such that \(\mathrm{H}_{\mathbf{M},\mathrm{F}'}^{(1)} \geqslant \mathrm{T}^{m}\mathrm{H}_{\mathbf{M},\mathbf{F}}^{(1)}\) .
\end{enumerate}

The part \(a)\) follows immediately from the definition of \(\mathbf{H}_{\mathbf{M},\mathbf{F}}^{(1)}\) ; the part \(b)\) follows from \(a)\) and Lemma 3.

THEOREM 2. --- Let A be a Noetherian ring, q an ideal of A, M a finitely generated A-module such that M/qM is nonzero and of finite length, and F a q-good filtration on M.

\begin{enumerate}
\def\labelenumi{\alph{enumi})}
\item
  There exists an integer \(d \geqslant 0\) and an element \(\mathbf{R}\) of \(\mathbf{Z}[\mathrm{T}, \mathrm{T}^{-1}]\) , uniquely determined, such that \(\mathbf{R}(1) > 0\) and \(\mathrm{H}_{\mathrm{M,F}} = (1 - \mathrm{T})^{-d}\mathbf{R}\) .
\item
  The integers d and R(1) are independent of the chosen q-good filtration F.
\item
  Consider the graded ring \(\operatorname{gr}(A)\) such that \(\operatorname{gr}_{n}(A) = \mathfrak{q}^{n}/\mathfrak{q}^{n+1}\) and the graded \(\operatorname{gr}(A)\)-module \(\operatorname{gr}(M)\) such that \(\operatorname{gr}_{n}(M) = F_{n}/F_{n+1}\). Since we have \(F_{n_{1}} = M\) and \(\mathfrak{q}F_{n} = F_{n+1}\) for \(n \geqslant n_{0}\), \(\operatorname{gr}(M)\) is generated by \(\bigoplus_{n_{1} \leqslant n \leqslant n_{0}} \operatorname{gr}_{n}(M)\), hence is of finite type. Moreover, if \((x_{1},\ldots,x_{r})\) is a finite generating system of the A-module \(\mathfrak{q}\), \(\operatorname{gr}(A)\) is generated by \(\operatorname{gr}_{0}(A)\) and the classes of the \(x_{i}\) modulo \(\mathfrak{q}^{2}\), hence is isomorphic to a graded quotient ring of \(H = (A/\mathfrak{q})[X_{1},\ldots,X_{r}]\). By th. 1 of no. 2, we have
\end{enumerate}

\[
(1 - \mathrm{T}) ^ {r} \mathrm{H} _ {\mathrm{M}, \mathrm{F}} \in \mathbf {Z} [ \mathrm{T}, \mathrm{T} ^ {- 1} ].
\]

We have \(H_{M,F} \neq 0\) and there therefore exist \(d \in N\) and \(R \in Z[T, T^{-1}]\) uniquely determined such that \(\mathrm{R}(1) > 0\) and \(\mathrm{H}_{\mathrm{M},\mathrm{F}} = (1 - \mathrm{T})^{-d}. R\) (Lemma 2 of no. 1).

\begin{enumerate}
\def\labelenumi{\alph{enumi})}
\setcounter{enumi}{1}
\tightlist
\item
  Let F' be another q-good filtration and write similarly
\end{enumerate}

\[
\mathrm{H} _ {\mathbf {M}, \mathrm{F} ^ {\prime}} = (1 - \mathrm{T}) ^ {- d ^ {\prime}} \mathbf {R} ^ {\prime}.
\]

By Prop. 3, b), there exists an integer m such that \((1 - T)^{-d'-1}R' \geqslant T^{m}(1 - T)^{-d-1}R\) . By Lemma 2, b) of no. 1, this implies \(d' \geqslant d\) and, if \(d' = d\) , \(R'(1) \geqslant R(1)\) . Exchanging the roles of F and \(F'\) , we obtain \(d = d'\) and \(R(1) = R'(1)\) .

Remark 1. --- With the notations of \(a\) ), write \(\mathbf{R} = \sum_{i \in \mathbf{Z}} a_i \mathbf{T}^i\) , and suppose \(d > 0\) . By no. 1, the relation \(\mathrm{H}_{\mathrm{M},\mathrm{F}} = (1 - \mathrm{T})^{-d}\mathrm{R}\) is also written

\[
\operatorname{long} _ {\mathbf {A}} \left(\mathrm{F} _ {n} / \mathrm{F} _ {n + 1}\right) = \sum_ {i \in \mathbf {Z}} a _ {i} \left[ \begin{array}{c} n - i + d - 1 \\ d - 1 \end{array} \right] = \sum_ {i \leqslant n} a _ {i} \binom {n - i + d - 1} {d - 1}.\tag{11}
\]

Similarly, since \(\mathbf{H}_{\mathbf{M},\mathbf{F}}^{(1)} = (1 - \mathbf{T})^{-d - 1}\mathbf{R}\) , one has

\[
\operatorname{long} _ {\mathbf {A}} \left(\mathrm{M} / \mathrm{F} _ {n + 1}\right) = \sum_ {i \in \mathbf {Z}} a _ {i} \left[ \begin{array}{c} n - i + d \\ d \end{array} \right] = \sum_ {i \leqslant n} a _ {i} \binom {n - i + d} {d}.\tag{12}
\]

Let A be a Noetherian ring, q an ideal of A, and M a finitely generated A-module such that M/qM has finite length. If M ≠ qM, then by Theorem 2, b), there exist integers \(d_{q}(M) \geqslant 0\) and \(e_{q}(M) > 0\) such that, for every q-good filtration F on M, there exists R ∈ Z{[}T, T \(^{-1}\) {]} with

\[
\mathrm{H} _ {\mathbf {M}, \mathrm{F}} = (1 - \mathrm{T}) ^ {- d _ {\mathfrak {q}} (\mathbf {M})} \mathrm{R}, \quad \mathrm{R} (1) = e _ {\mathfrak {q}} (\mathbf {M}).
\]

If M = qM, we set by convention \(d_{q}(M) = -\infty\) , \(e_{q}(M) = 0\) .

Remark 2. --- Saying that M/qM has finite length means that

\[
\operatorname{Supp} (\mathrm{M} / \mathfrak {q} \mathrm{M}) = \operatorname{Supp} (\mathrm{M}) \cap \mathrm{V} (\mathfrak {q})
\]

is formed of maximal ideals (IV, § 2, no. 5, prop. 7). We shall see below (no. 4, corollary to th. 3) that \(d_{\mathfrak{q}}(\mathbf{M})\) is the least upper bound of the numbers \(\dim_{\mathbf{A}_{\mathfrak{m}}}(M_{\mathfrak{m}})\) where m ranges over the set \(\operatorname{Supp}(\mathbf{M}) \cap \mathbf{V}(\mathfrak{q})\) .

COROLLARY. — Let A be a Noetherian ring, q an ideal of A, M a finitely generated A-module such that M/qM has finite length, and F a q-good filtration on M.

\begin{enumerate}
\def\labelenumi{\alph{enumi})}
\tightlist
\item
  For one to have \(d_{\mathfrak{q}}(\mathbf{M}) \leqslant 0\) it is necessary and sufficient that the sequence \((\mathfrak{q}^{n}\mathbf{M})\) be stationary, or equivalently that the sequence \((\mathrm{F}_{n})\) be stationary. Then, for all sufficiently large n,
\end{enumerate}

\[
\operatorname{long} \left(\mathrm{M} / \mathrm{F} _ {n + 1}\right) = \operatorname{long} \left(\mathrm{M} / \mathfrak {q} ^ {n + 1} \mathrm{M}\right) = e _ {\mathfrak {q}} (\mathrm{M}).
\]

\begin{enumerate}
\def\labelenumi{\alph{enumi})}
\setcounter{enumi}{1}
\tightlist
\item
  Suppose that we have \(d_{\mathfrak{q}}(\mathbf{M}) > 0\) . Then one has
\end{enumerate}

\begin{enumerate}
\def\labelenumi{(\arabic{enumi})}
\setcounter{enumi}{12}
\tightlist
\item
\end{enumerate}

\[
\operatorname{long} _ {\mathbf {A}} \left(\mathrm{F} _ {n} / \mathrm{F} _ {n + 1}\right) = e _ {\mathfrak {q}} (\mathrm{M}) n ^ {d _ {\mathfrak {q}} (\mathrm{M}) - 1} / \left(d _ {\mathfrak {q}} (\mathrm{M}) - 1\right)! + \rho_ {n} n ^ {d _ {\mathfrak {q}} (\mathrm{M}) - 2},\tag{14}
\]

\[
\operatorname{long} _ {\mathbf {A}} (\mathrm{M} / \mathrm{F} _ {n + 1}) = e _ {\mathfrak {q}} (\mathbf {M}) n ^ {d _ {\mathfrak {q}} (\mathbf {M})} / d _ {\mathfrak {q}} (\mathbf {M})! + \sigma_ {n} n ^ {d _ {\mathfrak {q}} (\mathbf {M}) - 1},
\]

where \(\rho_{n}\) and \(\sigma_{n}\) tend to a limit as \(n\) increases indefinitely.

This follows immediately from th. 2 and formula (2) of no. 1.

Remark 3. --- Suppose q is contained in the radical of A. Then, by Nakayama's lemma, the sequence (q \(^{n}\) M) is stationary if and only if one has q \(^{n}\) M = 0 for n sufficiently large. It then follows from part a) of the corollary that one has d \(_{q}\) (M) ≤ 0 if and only if M is of finite length, and one then has e \(_{q}\) (M) = long \(_{A}\) (M).

PROPOSITION 4. — Let A be a Noetherian ring, \(x_{1},...,x_{r}\) elements of A, \(\mathfrak{x}\) the ideal they generate, and M a finitely generated A-module such that \(M/\mathfrak{x}M\) is nonzero and of finite length.

\begin{enumerate}
\def\labelenumi{\alph{enumi})}
\item
  We have \(d_{\mathfrak{x}}(\mathbf{M}) \leqslant r\) .
\item
  If \(d_{\mathfrak{x}}(M) = r\), then \(e_{\mathfrak{x}}(M) \leqslant \text{long}_{A}(M/\mathfrak{x}M)\).
\item
  If the sequence \((x_{1},...,x_{r})\) is completely secant for M (A, X, p. 157), then \(d_{\mathfrak{x}}(\mathbf{M}) = r\) and \(e_{\mathfrak{x}}(\mathbf{M}) = \text{long}_{\mathbf{A}}(\mathbf{M}/\mathfrak{x}\mathbf{M})\). The converse is true if the \(x_{i}\) belong to the radical of A.
\end{enumerate}

Let H be the polynomial ring \((\mathbf{A} / \mathfrak{x})[\mathbf{X}_1,\dots ,\mathbf{X}_r]\). We equip \(\mathbf{G} = \bigoplus_{n}\mathfrak{x}^{n}\mathbf{M} / \mathfrak{x}^{n + 1}\mathbf{M}\)

of the structure of graded H-module for which \((\mathrm{X}_{i})_{\mathrm{G}}\) is multiplication by the class of \(x_{i}\) modulo \(\mathfrak{x}^{2}\). With the notations \(P_{G}, Q_{G}, c_{G}\) from No. 2, we have \(H_{M,\mathfrak{x}} = P_{G}\) , hence \((1 - T)^{-d_{\mathfrak{x}}(M)}R = (1 - T)^{-r}Q_{G}\) , where \(R(1) = e_{\mathfrak{x}}(M) > 0\) and \(Q_{G}(1) = c_{G}\) .

Thus either \(d_{\mathfrak{x}}(\mathbf{M}) < r\) and \(c_{\mathbf{G}} = 0\), or \(d_{\mathfrak{x}}(\mathbf{M}) = r\) and \(c_{\mathbf{G}} = e_{\mathfrak{x}}(\mathbf{M})\). Moreover, by prop. 1 of no. 2, we have \(c_{\mathbf{G}} \leqslant \text{long}(\mathbf{M}/\mathfrak{x}\mathbf{M})\), and there is equality if and only if the canonical homomorphism \(\mathrm{A}/\mathfrak{x}[\mathrm{X}_{1}, ..., \mathrm{X}_{r}] \otimes_{\mathrm{A}/\mathfrak{x}} \mathrm{M}/\mathfrak{x}\mathrm{M} \to \bigoplus_{n}\mathfrak{x}^{n}\mathrm{M}/\mathfrak{x}^{n+1}\mathrm{M}\) is bijective.

This entails the proposition, in view of A, X, p. 160, th. 1.

PROPOSITION 5. --- Let \(0 \to \mathbf{M}' \to \mathbf{M} \to \mathbf{M}'' \to 0\) be an exact sequence of finitely generated modules over a Noetherian ring A, and q an ideal of A.

\begin{enumerate}
\def\labelenumi{\alph{enumi})}
\item
  For \(M/\mathfrak{q}M\) to have finite length, it is necessary and sufficient that the same be true of \(M'/\mathfrak{q}M'\) and \(M''/\mathfrak{q}M''\).
\item
  Suppose M/qM has finite length. Then we are in one of the following three cases:
\end{enumerate}

\[
1) d _ {\mathfrak {q}} (\mathbf {M}) = d _ {\mathfrak {q}} \left(\mathbf {M} ^ {\prime}\right) > d _ {\mathfrak {q}} \left(\mathbf {M} ^ {\prime \prime}\right) and e _ {\mathfrak {q}} (\mathbf {M}) = e _ {\mathfrak {q}} \left(\mathbf {M} ^ {\prime}\right),
\]

\begin{enumerate}
\def\labelenumi{\arabic{enumi})}
\setcounter{enumi}{1}
\tightlist
\item
  \(d_{\mathfrak{q}}(\mathbf{M}) = d_{\mathfrak{q}}(\mathbf{M}^{\prime \prime}) > d_{\mathfrak{q}}(\mathbf{M}^{\prime})\) and \(e_{\mathfrak{q}}(\mathbf{M}) = e_{\mathfrak{q}}(\mathbf{M}^{\prime \prime}),\)
\end{enumerate}

\[
3) d _ {\mathfrak {q}} (\mathbf {M}) = d _ {\mathfrak {q}} \left(\mathbf {M} ^ {\prime}\right) = d _ {\mathfrak {q}} \left(\mathbf {M} ^ {\prime \prime}\right) and e _ {\mathfrak {q}} (\mathbf {M}) = e _ {\mathfrak {q}} \left(\mathbf {M} ^ {\prime}\right) + e _ {\mathfrak {q}} \left(\mathbf {M} ^ {\prime \prime}\right).
\]

\begin{enumerate}
\def\labelenumi{\alph{enumi})}
\item
  We have \(\operatorname{Supp}(\mathbf{M}) = \operatorname{Supp}(\mathbf{M}') \cup \operatorname{Supp}(\mathbf{M}'')\) and the assertion follows from Remark 2.
\item
  Equip M with a q-good filtration F (for example, the q-adic filtration), M'' with the quotient filtration F'', and M' with the induced filtration F'. The filtrations F' and F'' are q-good (III, § 3, no. 1, prop. 1). Then for each n we have an exact sequence of A-modules
\end{enumerate}

\[
0 \rightarrow \mathrm{F} _ {n} ^ {\prime} / \mathrm{F} _ {n + 1} ^ {\prime} \rightarrow \mathrm{F} _ {n} / \mathrm{F} _ {n + 1} \rightarrow \mathrm{F} _ {n} ^ {\prime \prime} / \mathrm{F} _ {n + 1} ^ {\prime \prime} \rightarrow 0
\]

(III, § 2, no. 4, prop. 2), so that we have \(\mathrm{H}_{\mathbf{M},\mathbf{F}} = \mathrm{H}_{\mathbf{M}',\mathbf{F}'} + \mathrm{H}_{\mathbf{M}'',\mathbf{F}''}\) , or equivalently

\[
(1 - \mathrm{T}) ^ {- d _ {\mathfrak {q}} (\mathrm{M})} \mathrm{R} = (1 - \mathrm{T}) ^ {- d _ {\mathfrak {q}} (\mathrm{M} ^ {\prime})} \mathrm{R} ^ {\prime} + (1 - \mathrm{T}) ^ {- d _ {\mathfrak {q}} (\mathrm{M} ^ {\prime \prime})} \mathrm{R} ^ {\prime \prime},
\]

with \(R,R',R''\in\mathbf{Z}[T,T^{-1}]\), \(R(1)=e_{\mathfrak{q}}(\mathbf{M})\), \(R'(1)=e_{\mathfrak{q}}(\mathbf{M}')\), \(R''(1)=e_{\mathfrak{q}}(\mathbf{M}'')\). Assertion b) follows immediately from this.

